% Figure source fragment: schematic hyperbolic-tangent saturation curve.
% This file is inserted inside a figure environment using \input.
% The curve is schematic and is not a plot of experimental measurements.
\begin{center}
\setlength{\unitlength}{0.68cm}
\begin{picture}(8.4,4.4)
    % Coordinate axes.
    \put(0.45,2.15){\vector(1,0){7.45}}
    \put(4.10,0.30){\vector(0,1){3.80}}

    % Approximate S-shaped curve using quadratic Bezier segments.
    \qbezier(0.75,0.55)(1.55,0.55)(2.35,0.92)
    \qbezier(2.35,0.92)(3.10,1.38)(3.55,1.83)
    \qbezier(3.55,1.83)(3.83,2.08)(4.10,2.15)
    \qbezier(4.10,2.15)(4.37,2.22)(4.65,2.47)
    \qbezier(4.65,2.47)(5.10,2.92)(5.85,3.38)
    \qbezier(5.85,3.38)(6.65,3.75)(7.45,3.75)

    % Saturation guides.
    \put(0.75,3.75){\line(1,0){6.70}}
    \put(0.75,0.55){\line(1,0){6.70}}

    % Axis labels and function annotation.
    \put(7.95,1.93){\(x\)}
    \put(4.23,4.02){\(y\)}
    \put(7.10,3.86){\(1\)}
    \put(7.10,0.20){\(-1\)}
    \put(4.22,1.88){\(0\)}
    \put(1.05,2.38){\(\tanh(\alpha x)\)}
\end{picture}
\end{center}
